\chapter{线性神经网络}

线性模型把输入 $\bm{x}$ 映成标量或类别分数，再用平方损失或交叉熵得到可微目标。
训练在小批量上用梯度更新参数；分类时在线性输出后再接 $\mathrm{softmax}$。
下面按原书顺序写出对象、损失、最优条件和实现所对应的映射。

\section{线性回归}

回归估计 $y\in\mathbb{R}$。给定特征 $\bm{x}\in\mathbb{R}^{d}$，模型输出一个实数预测 $\hat{y}$。

\subsection{线性回归的基本元素}

线性回归假设 $\hat{y}$ 对 $\bm{x}$ 仿射，噪声为加性高斯。训练集
\[
\mathcal{D}=\{(\bm{x}^{(i)},y^{(i)})\}_{i=1}^{n}
\]
用来确定权重 $\bm{w}$ 与偏置 $b$。

\subsubsection{线性模型}

房价例子把面积与房龄写成加权和
\[
\mathrm{price}=w_{\mathrm{area}}\cdot\mathrm{area}+w_{\mathrm{age}}\cdot\mathrm{age}+b.
\]
一般地，$d$ 个特征给出
\[
\hat{y}=w_1 x_1+\cdots+w_d x_d+b=\bm{w}^{\top}\bm{x}+b.
\]
把 $n$ 个样本排成 $\bm{X}\in\mathbb{R}^{n\times d}$ 后，向量形式为
\[
\hat{\bm{y}}=\bm{X}\bm{w}+b.
\]

\subsubsection{损失函数}

第 $i$ 个样本的平方损失为
\[
l^{(i)}(\bm{w},b)=\frac{1}{2}\bigl(\hat{y}^{(i)}-y^{(i)}\bigr)^{2}.
\]
经验风险是样本平均
\begin{align*}
L(\bm{w},b)
&=\frac{1}{n}\sum_{i=1}^{n}l^{(i)}(\bm{w},b)\\
&=\frac{1}{n}\sum_{i=1}^{n}\frac{1}{2}\bigl(\bm{w}^{\top}\bm{x}^{(i)}+b-y^{(i)}\bigr)^{2}.
\end{align*}
参数估计定义为
\[
\bm{w}^{\ast},b^{\ast}=\operatorname*{argmin}_{\bm{w},b}\,L(\bm{w},b).
\]

\subsubsection{解析解}

把偏置并入 $\bm{X}$ 的全 $1$ 列后，正规方程给出
\[
\bm{w}^{\ast}=(\bm{X}^{\top}\bm{X})^{-1}\bm{X}^{\top}\bm{y},
\]
成立条件是 $\bm{X}^{\top}\bm{X}$ 可逆，即特征线性无关且 $n$ 足够大。

\subsubsection{随机梯度下降}

小批量 $\mathcal{B}$ 上的更新为
\[
(\bm{w},b)
\leftarrow
(\bm{w},b)
-
\frac{\eta}{|\mathcal{B}|}
\sum_{i\in\mathcal{B}}
\partial_{(\bm{w},b)}l^{(i)}(\bm{w},b).
\]
平方损失的显式梯度是
\begin{align*}
\bm{w}
&\leftarrow
\bm{w}
-
\frac{\eta}{|\mathcal{B}|}
\sum_{i\in\mathcal{B}}
\bm{x}^{(i)}
\bigl(\bm{w}^{\top}\bm{x}^{(i)}+b-y^{(i)}\bigr),\\
b
&\leftarrow
b
-
\frac{\eta}{|\mathcal{B}|}
\sum_{i\in\mathcal{B}}
\bigl(\bm{w}^{\top}\bm{x}^{(i)}+b-y^{(i)}\bigr).
\end{align*}

\subsubsection{用模型进行预测}

学得 $(\hat{\bm{w}},\hat{b})$ 后，新样本的预测是仿射映射
\[
\hat{y}=\hat{\bm{w}}^{\top}\bm{x}+\hat{b}.
\]

\subsection{矢量化加速}

对 $\bm{a},\bm{b}\in\mathbb{R}^{m}$，逐元素求和应写成一次向量加法
\[
\bm{c}=\bm{a}+\bm{b},\qquad c_{i}=a_{i}+b_{i},
\]
而不是 $m$ 次标量循环。小批量上的 $\bm{X}\bm{w}$ 同样一次完成，复杂度由线性代数库承担。

\subsection{正态分布与平方损失}

高斯密度为
\[
p(x)=\frac{1}{\sqrt{2\pi\sigma^{2}}}
\exp\Bigl(-\frac{1}{2\sigma^{2}}(x-\mu)^{2}\Bigr).
\]
观测模型
\[
y=\bm{w}^{\top}\bm{x}+b+\epsilon
\]
在 $\epsilon\sim\mathcal{N}(0,\sigma^{2})$ 下给出
\[
P(y\mid\bm{x})
=
\frac{1}{\sqrt{2\pi\sigma^{2}}}
\exp\Bigl(-\frac{1}{2\sigma^{2}}(y-\bm{w}^{\top}\bm{x}-b)^{2}\Bigr).
\]
独立样本的似然是
\[
P(\bm{y}\mid\bm{X})=\prod_{i=1}^{n}p(y^{(i)}\mid\bm{x}^{(i)}),
\]
负对数似然为
\begin{align*}
-\log P(\bm{y}\mid\bm{X})
&=
\sum_{i=1}^{n}\frac{1}{2}\log(2\pi\sigma^{2})\\
&\qquad
+\frac{1}{2\sigma^{2}}\bigl(y^{(i)}-\bm{w}^{\top}\bm{x}^{(i)}-b\bigr)^{2}.
\end{align*}
$\sigma$ 固定时，极大似然与最小化 $L(\bm{w},b)$ 等价。

\subsection{从线性回归到深度网络}

同一仿射映射可看成只有一层计算单元的网络：输入给定，输出由 $\bm{w},b$ 计算。

\subsubsection{神经网络图}

输入层提供 $x_{1},\ldots,x_{d}$，输出层一个单元
\[
o_{1}=\bm{w}^{\top}\bm{x}+b.
\]
每个输入都连到该单元，参数个数为 $d+1$。

\subsubsection{生物学}

树突接收 $x_{i}$，突触权重为 $w_{i}$，胞体做加权和后再经激活
\[
o=\phi\bigl(\bm{w}^{\top}\bm{x}+b\bigr).
\]
线性回归对应 $\phi$ 为恒等映射。

\subsection{小结}

对象是仿射预测 $\hat{y}=\bm{w}^{\top}\bm{x}+b$，目标是 $L(\bm{w},b)$，
高斯噪声下它等于负对数似然；可逆时用正规方程，否则用小批量梯度。

\subsection{练习}

核验 $\sum_{i}(x_{i}-b)^{2}$ 对 $b$ 的驻点、把 $b$ 并入 $\bm{X}$ 后正规方程是否仍成立，
以及 $\bm{X}^{\top}\bm{X}$ 奇异时梯度法与解析解的关系。

\section{线性回归的从零开始实现}

实现对应四个映射：采样数据、按批读取、仿射预测、平方损失上的小批量更新。

\subsection{生成数据集}

真值为
\[
\bm{y}=\bm{X}\bm{w}+b+\bm{\epsilon},
\]
其中 $\bm{\epsilon}$ 的分量独立同分布。合成数据时 $\bm{w},b$ 已知，用来对照估计。

\subsection{读取数据集}

一次迭代取出小批量
\[
(\bm{X}_{\mathcal{B}},\bm{y}_{\mathcal{B}}),\qquad
|\mathcal{B}|=B.
\]
遍历前打乱下标，使各批近似无偏。

\subsection{初始化模型参数}

\[
\bm{w}\sim\mathcal{N}(\bm{0},0.01^{2}\bm{I}),\qquad b=0,
\]
并对其后每次损失关于 $(\bm{w},b)$ 求梯度。

\subsection{定义模型}

前向是矩阵--向量积加偏置广播
\[
\hat{\bm{y}}=\bm{X}\bm{w}+b.
\]

\subsection{定义损失函数}

批量平方损失与上一节一致，实现时把 $\bm{y}$ 的形状对齐到 $\hat{\bm{y}}$：
\[
l=\frac{1}{2}\|\hat{\bm{y}}-\bm{y}\|_{2}^{2}
\quad\text{（再按约定取平均或求和）}.
\]

\subsection{定义优化算法}

一步小批量梯度下降
\[
\bm{\theta}
\leftarrow
\bm{\theta}
-
\frac{\eta}{|\mathcal{B}|}\nabla_{\bm{\theta}}
\sum_{i\in\mathcal{B}}l^{(i)},
\qquad
\bm{\theta}=(\bm{w},b).
\]
用 $|\mathcal{B}|$ 规范化后，步长不随批量大小线性膨胀。

\subsection{训练}

每个周期对全部小批量执行
\begin{align*}
\hat{\bm{y}}&\leftarrow \bm{X}_{\mathcal{B}}\bm{w}+b,\\
g&\leftarrow
\partial_{(\bm{w},b)}\frac{1}{|\mathcal{B}|}\sum_{i\in\mathcal{B}}l^{(i)},\\
(\bm{w},b)&\leftarrow(\bm{w},b)-\eta g.
\end{align*}

\subsection{小结}

从零实现把训练写成张量上的 $\hat{\bm{y}}=\bm{X}\bm{w}+b$ 与
$(\bm{w},b)\leftarrow(\bm{w},b)-\eta g$，不引入层对象。

\subsection{练习}

核验 $\bm{w}=\bm{0}$ 时梯度是否仍指向残差方向、
二阶导数在自动微分中的形状，以及 $n$ 不能整除 $|\mathcal{B}|$ 时最后一批的基数。

\section{线性回归的简洁实现}

同一映射改由层、损失类与优化器组合，前向与更新规则不变。

\subsection{生成数据集}

仍用 $\bm{y}=\bm{X}\bm{w}+b+\bm{\epsilon}$ 生成 $(\bm{X},\bm{y})$。

\subsection{读取数据集}

数据迭代器实现
\[
(\bm{X},\bm{y})
\;\longmapsto\;
\{(\bm{X}_{\mathcal{B}},\bm{y}_{\mathcal{B}})\}.
\]
训练模式下每个周期打乱样本次序。

\subsection{定义模型}

全连接层
\[
f(\bm{X})=\bm{X}\bm{W}+\bm{b},\qquad \bm{W}\in\mathbb{R}^{d\times 1}
\]
作为唯一计算层；顺序容器只含这一层。

\subsection{初始化模型参数}

\[
W_{ij}\sim\mathcal{N}(0,0.01^{2}),\qquad b=0.
\]

\subsection{定义损失函数}

均方误差是平方 $L_{2}$ 的样本平均
\[
L=\frac{1}{n}\sum_{i=1}^{n}\frac{1}{2}\bigl(\hat{y}^{(i)}-y^{(i)}\bigr)^{2}.
\]

\subsection{定义优化算法}

优化器在参数集合上执行
\[
\bm{\theta}\leftarrow\bm{\theta}-\eta\nabla_{\bm{\theta}}L_{\mathcal{B}}.
\]

\subsection{训练}

每个小批量依次：前向 $L_{\mathcal{B}}$、反向 $\nabla_{\bm{\theta}}L_{\mathcal{B}}$、参数更新。
周期结束时用 $L$ 监控拟合。

\subsection{小结}

简洁实现把 $\bm{X}\mapsto\bm{X}\bm{W}+\bm{b}$、平方损失与小批量梯度封装成可组合对象，数学映射与从零实现相同。

\subsection{练习}

核验分段损失
\[
l(y,y')=
\begin{cases}
|y-y'|-\dfrac{\sigma}{2}, & |y-y'|>\sigma,\\[0.4em]
\dfrac{1}{2\sigma}(y-y')^{2}, & \text{其它情况}
\end{cases}
\]
在 $|y-y'|=\sigma$ 处的连续性与可导性。

\section{softmax回归}

分类不估计“多少”，而估计“哪一类”。线性层输出 $q$ 个分数，再用 $\mathrm{softmax}$ 变成单纯形上的概率。

\subsection{分类问题}

$q=3$ 时标签用独热向量
\[
y\in\{(1,0,0),\ (0,1,0),\ (0,0,1)\}.
\]
一般 $y\in\{1,\ldots,q\}$ 对应 $\bm{y}=\bm{e}_{y}\in\{0,1\}^{q}$。

\subsection{网络架构}

四个特征、三个类别时，线性分数为
\begin{align*}
o_{1}&=x_{1}w_{11}+x_{2}w_{12}+x_{3}w_{13}+x_{4}w_{14}+b_{1},\\
o_{2}&=x_{1}w_{21}+x_{2}w_{22}+x_{3}w_{23}+x_{4}w_{24}+b_{2},\\
o_{3}&=x_{1}w_{31}+x_{2}w_{32}+x_{3}w_{33}+x_{4}w_{34}+b_{3}.
\end{align*}
矩阵形式 $\bm{o}=\bm{W}\bm{x}+\bm{b}$，$\bm{W}\in\mathbb{R}^{q\times d}$。

\subsection{全连接层的参数开销}

$d$ 入 $q$ 出的全连接层参数个数为
\[
\#(\bm{W},\bm{b})=dq+q=\mathcal{O}(dq).
\]

\subsection{softmax运算}

把分数映到概率单纯形
\[
\hat{\bm{y}}=\mathrm{softmax}(\bm{o}),
\qquad
\hat{y}_{j}=\frac{\exp(o_{j})}{\sum_{k}\exp(o_{k})}.
\]
指数严格递增，故
\[
\operatorname*{argmax}_{j}\hat{y}_{j}=\operatorname*{argmax}_{j}o_{j}.
\]

\subsection{小批量样本的矢量化}

$\bm{X}\in\mathbb{R}^{n\times d}$ 时
\begin{align*}
\bm{O}&=\bm{X}\bm{W}+\bm{b},\\
\hat{\bm{Y}}&=\mathrm{softmax}(\bm{O}),
\end{align*}
其中 $\mathrm{softmax}$ 沿类别轴作用。

\subsection{损失函数}

仍由最大似然导出：独立样本下最小化 $-\log P(\bm{Y}\mid\bm{X})$。

\subsubsection{对数似然}

\[
P(\bm{Y}\mid\bm{X})=\prod_{i=1}^{n}P(\bm{y}^{(i)}\mid\bm{x}^{(i)}).
\]
负对数似然分解为交叉熵之和
\begin{align*}
-\log P(\bm{Y}\mid\bm{X})
&=
\sum_{i=1}^{n}-\log P(\bm{y}^{(i)}\mid\bm{x}^{(i)})
=
\sum_{i=1}^{n}l(\bm{y}^{(i)},\hat{\bm{y}}^{(i)}),
\end{align*}
其中
\[
l(\bm{y},\hat{\bm{y}})=-\sum_{j=1}^{q}y_{j}\log\hat{y}_{j}.
\]
独热标签时 $l=-\log\hat{y}_{y^{\ast}}$。

\subsubsection{softmax及其导数}

把 $\hat{y}_{j}$ 代入交叉熵：
\begin{align*}
l(\bm{y},\hat{\bm{y}})
&=
-\sum_{j=1}^{q}y_{j}\log\frac{\exp(o_{j})}{\sum_{k=1}^{q}\exp(o_{k})}\\
&=
\sum_{j=1}^{q}y_{j}\log\sum_{k=1}^{q}\exp(o_{k})
-\sum_{j=1}^{q}y_{j}o_{j}\\
&=
\log\sum_{k=1}^{q}\exp(o_{k})
-\sum_{j=1}^{q}y_{j}o_{j}.
\end{align*}
对分数求导得到
\begin{align*}
\partial_{o_{j}}l(\bm{y},\hat{\bm{y}})
&=
\frac{\exp(o_{j})}{\sum_{k=1}^{q}\exp(o_{k})}-y_{j}\\
&=
\mathrm{softmax}(\bm{o})_{j}-y_{j}.
\end{align*}
因此反向传播中，$\nabla_{\bm{o}}l=\hat{\bm{y}}-\bm{y}$。

\subsubsection{交叉熵损失}

若 $\bm{y}$ 本身已是分布而非独热，同一式
\[
l(\bm{y},\hat{\bm{y}})=-\sum_{j=1}^{q}y_{j}\log\hat{y}_{j}
\]
仍是交叉熵，等于在标签分布下对 $-\log\hat{y}_{j}$ 的期望。

\subsection{信息论基础}

交叉熵同时度量编码长度与似然。

\subsubsection{熵}

分布 $P$ 的熵
\[
H[P]=\sum_{j}-P(j)\log P(j)
\]
是按 $P$ 编码时的最小期望码长。

\subsubsection{信息量}

事件 $j$ 的自信息为
\[
I(j)=-\log P(j)=\log\frac{1}{P(j)}.
\]
$P(j)$ 越小，$I(j)$ 越大。熵是自信息的期望。

\subsubsection{重新审视交叉熵}

从 $P$ 到 $Q$ 的交叉熵
\[
H(P,Q)=\sum_{j}-P(j)\log Q(j)=H(P)+D_{\mathrm{KL}}(P\|Q)
\]
在 $Q=P$ 时取最小 $H(P)$。分类目标等价于极大似然，也等价于降低用 $\hat{\bm{y}}$ 编码真实标签所需的惊异。

\subsection{模型预测和评估}

预测类别
\[
\hat{y}=\operatorname*{argmax}_{j}\hat{y}_{j}.
\]
精度是正确个数与总数之比
\[
\mathrm{acc}=\frac{1}{n}\sum_{i=1}^{n}\mathbf{1}\{\hat{y}^{(i)}=y^{(i)}\}.
\]
精度不可导，训练仍最小化交叉熵。

\subsection{小结}

$\mathrm{softmax}$ 把 $\bm{o}$ 映到概率，$\nabla_{\bm{o}}l=\hat{\bm{y}}-\bm{y}$ 连接交叉熵与线性层；
精度在单纯形顶点上评估 $\operatorname*{argmax}$。

\subsection{练习}

核验交叉熵对 $\bm{o}$ 的 Hessian 与 $\mathrm{softmax}(\bm{o})$ 的协方差是否一致，
以及均匀分布 $(\tfrac{1}{3},\tfrac{1}{3},\tfrac{1}{3})$ 下二元码的冗余。

\section{图像分类数据集}

Fashion-MNIST 提供 $q=10$ 的服装灰度图，后续 softmax 实验都在该 $\mathcal{D}$ 上评估。

\subsection{读取数据集}

训练集 $n_{\mathrm{train}}=60000$，测试集 $n_{\mathrm{test}}=10000$，
每类分别 $6000$ 与 $1000$ 张。单张图
\[
\bm{X}\in\mathbb{R}^{28\times 28},\qquad C=1,
\]
形状记为 $h\times w=28\times 28$。测试集不参与参数更新。

\subsection{读取小批量}

迭代器输出
\[
\mathsf{X}\in\mathbb{R}^{|\mathcal{B}|\times 1\times 28\times 28},\qquad
\bm{y}\in\{1,\ldots,10\}^{|\mathcal{B}|}.
\]
训练时打乱，测试时按序。

\subsection{整合所有组件}

读取映射把 $(\text{批量大小},\text{可选边长})$ 映到训练与测试迭代器。
可选边长把 $28\times 28$ 双线性映到指定 $(h',w')$。

\subsection{小结}

图像对象是 $h\times w$ 灰度阵，批量读取把它们变成可并行的 $(\mathsf{X},\bm{y})$。

\subsection{练习}

核验 $|\mathcal{B}|\to 1$ 时 I/O 相对计算的占比，以及迭代器是否成为训练的时间瓶颈。

\section{softmax回归的从零开始实现}

把 $28\times 28$ 展平为 $\mathbb{R}^{784}$，线性层输出 $10$ 维，再接 $\mathrm{softmax}$ 与交叉熵。

\subsection{初始化模型参数}

\[
\bm{W}\in\mathbb{R}^{784\times 10},\qquad \bm{b}\in\mathbb{R}^{1\times 10},
\]
$W_{ij}\sim\mathcal{N}(0,\sigma^{2})$，$b_{j}=0$。展平是
\[
\mathbb{R}^{28\times 28}\longrightarrow\mathbb{R}^{784}.
\]

\subsection{定义softmax操作}

对矩阵按行归一化
\[
\mathrm{softmax}(\bm{X})_{ij}=\frac{\exp(X_{ij})}{\sum_{k}\exp(X_{ik})}.
\]

\subsection{定义模型}

\[
\hat{\bm{Y}}=\mathrm{softmax}(\bm{X}\bm{W}+\bm{b}),
\]
其中 $\bm{X}$ 已展平为 $|\mathcal{B}|\times 784$。

\subsection{定义损失函数}

独热选取使
\[
l(\bm{y},\hat{\bm{Y}})=-\sum_{i=1}^{|\mathcal{B}|}\log\hat{Y}_{i,y^{(i)}}.
\]

\subsection{分类精度}

\[
\mathrm{acc}(\hat{\bm{Y}},\bm{y})
=
\frac{1}{|\mathcal{B}|}
\sum_{i}
\mathbf{1}\bigl\{\operatorname*{argmax}_{j}\hat{Y}_{ij}=y^{(i)}\bigr\}.
\]

\subsection{训练}

与线性回归同一循环：前向 $\hat{\bm{Y}}$、交叉熵、反向、小批量更新。
每个周期在测试集上报告 $\mathrm{acc}$。

\subsection{预测}

对新图取
\[
\hat{y}=\operatorname*{argmax}_{j}\mathrm{softmax}(\bm{x}^{\top}\bm{W}+\bm{b})_{j}.
\]

\subsection{小结}

从零实现的核心是
$\hat{\bm{Y}}=\mathrm{softmax}(\bm{X}\bm{W}+\bm{b})$ 与 $l=-\sum\log\hat{y}_{y}$，
训练循环与线性回归同构。

\subsection{练习}

核验 $\exp(o_{j})$ 在 $o_{j}$ 很大时的溢出、$\log\hat{y}_{j}$ 在 $\hat{y}_{j}\to 0$ 时的定义域，
以及 $\operatorname*{argmax}$ 是否总是决策最优。

\section{softmax回归的简洁实现}

同一分类映射改由全连接层加交叉熵实现；数值上把 $\mathrm{softmax}$ 与 $\log$ 合并。

\subsection{初始化模型参数}

顺序模型中一层 $\mathbb{R}^{784}\to\mathbb{R}^{10}$，
\[
W_{ij}\sim\mathcal{N}(0,0.01^{2}),\qquad b_{j}=0.
\]

\subsection{重新审视Softmax的实现}

平移不变性给出数值稳定形式
\begin{align*}
\hat{y}_{j}
&=
\frac{\exp(o_{j}-\max_{k}o_{k})\exp(\max_{k}o_{k})}
{\sum_{k}\exp(o_{k}-\max_{k}o_{k})\exp(\max_{k}o_{k})}\\
&=
\frac{\exp(o_{j}-\max_{k}o_{k})}{\sum_{k}\exp(o_{k}-\max_{k}o_{k})}.
\end{align*}
对数概率为
\begin{align*}
\log\hat{y}_{j}
&=
\log\frac{\exp(o_{j}-\max_{k}o_{k})}{\sum_{k}\exp(o_{k}-\max_{k}o_{k})}\\
&=
o_{j}-\max_{k}o_{k}
-\log\sum_{k}\exp(o_{k}-\max_{k}o_{k}).
\end{align*}
交叉熵直接用 $\log\hat{y}_{y^{\ast}}$，避免先指数再取对数。

\subsection{优化算法}

小批量随机梯度
\[
\bm{\theta}\leftarrow\bm{\theta}-\eta\nabla_{\bm{\theta}}L_{\mathcal{B}},
\qquad \eta=0.1.
\]

\subsection{训练}

调用与从零实现相同的周期循环；框架在交叉熵内部完成稳定的 $\log\mathrm{softmax}$。

\subsection{小结}

简洁路径仍是 $\bm{O}=\bm{X}\bm{W}+\bm{b}$ 与交叉熵，差别在于用
$\log\hat{y}_{j}=o_{j}-\max o-\mathrm{LSE}(\bm{o}-\max o)$ 保证数值稳定。

\subsection{练习}

核验增大周期后测试精度下降是否对应过拟合，以及 $|\mathcal{B}|,\eta$ 对 $\nabla L_{\mathcal{B}}$ 方差的影响。
